Fig.~\ref{fig:interp} looks inside the plant-inference attention. The attention of the context token over the window is nearly uniform, between \interpAttMin{} and \interpAttMax{} against $1/L=\interpAttUniform{}$ (a): the token pools the whole window rather than a few samples, and the plant information sits in the pooled physics features. Panels (b) and (c) follow the $R^2$ between the context-token estimate and the true parameter along the transition on the unseen plants: each parameter is resolved while the response that reveals it is inside the window and released once the window has slid past it. The lift-curve offset is read off the wing-borne trim at the start of the back transition ($R^2=\interpRtwoCLzero{}$ at \interpTCLzero{}~s), the rotor gain once the rotors respond in the \interpTrBetaR{} transition ($\interpRtwoBetaR{}$ at \interpTBetaR{}~s), the attitude-loop rate from the first pitch response in the \interpTrKappa{} transition ($\interpRtwoKappa{}$ at \interpTKappa{}~s) and the lift-curve slope during the acceleration ($\interpRtwoCLa{}$ at \interpTCLa{}~s); the drag parameters, which act on the airspeed slowly, are not resolved ($R^2$ at most \interpRtwoCDzero{}). The estimates are shrunk towards the centre of the box by the measurement noise, to \interpShrinkCLzero{} of the true spread for $C_{L0}$, as a least-squares head under noise should be, so (d) and (e) show them standardised. A sliding window implies this forgetting; the ablation of the auxiliary head (Section~\ref{sec:ablation}) measures what the estimate is worth to the policy.
\begin{figure*}[t]\centering\includegraphics[width=\textwidth]{fig13_interp.pdf}
\caption{Interpretability on the unseen plants: (a) attention of the context token over the window, front and back, against the uniform weight $1/L$; (b), (c) $R^2$ of the context-token estimate of the lift-curve offset, the lift-curve slope, the attitude-loop rate and the rotor gain against the true values along the front and the back transition; (d), (e) standardised estimates against the standardised true values at the time of the best $R^2$.}\label{fig:interp}
\end{figure*}
